\documentclass[11pt]{article}
\usepackage[T1]{fontenc}
\usepackage[margin=1in]{geometry}
\usepackage{amsmath,amssymb,amsthm,mathtools,hyperref}
\hypersetup{hidelinks}
\newtheorem{theorem}{Theorem}[section]
\newtheorem{lemma}[theorem]{Lemma}
\newtheorem{proposition}[theorem]{Proposition}
\newtheorem{corollary}[theorem]{Corollary}
\theoremstyle{remark}\newtheorem{remark}[theorem]{Remark}
\newcommand{\Q}{\mathbf Q}\newcommand{\C}{\mathbf C}
\newcommand{\Z}{\mathbf Z}\newcommand{\F}{\mathbf F}
\newcommand{\CH}{\operatorname{CH}}\newcommand{\Sym}{\operatorname{Sym}}
\newcommand{\im}{\operatorname{im}}\newcommand{\id}{\operatorname{id}}
\newcommand{\cl}{\operatorname{cl}}\newcommand{\Gal}{\operatorname{Gal}}
\newcommand{\Alt}{\operatorname{Alt}}
\title{Determinants of cyclic covers of the line:\\a finite-quotient reduction to Fermat varieties}
\author{Kyle Kistner}
\date{5 September 2026 --- proof and attribution audit draft}
\begin{document}\maketitle
\begin{abstract}
Let $C\to\mathbf P^1$ be a cyclic cover of prime degree $p$ with $s\geq3$
branch points. Put $h=s-2$. We exhibit the rational cyclotomic determinant
of $H^1(C)$ as an algebraic direct summand of the middle cohomology of a
degree-$p$ Fermat hypersurface of dimension $h$. The construction uses two
finite maps to the same normal quotient and an explicit Chow-correspondence
identity; it requires no specialization or singular-cohomology comparison.
Shioda's published theorem for arbitrary products of prime-degree Fermat
varieties then implies algebraicity of every rational Hodge class in every
tensor product of these determinant spaces. In particular it implies the
balanced mixed-determinant statement for arbitrary lists of cyclic covers.
This supplies an attribution-sensitive simplification of the companion
fusion argument. The Fermat-quotient construction has direct precedent in
Schoen's 1989 paper; no priority claim is made for the construction or its
product consequence. This is a classical proof draft, not a formalization
or an independent referee certification.
\end{abstract}

\section{Statement and exact rational realization}
All varieties and correspondences are over $\C$, and Chow groups have
rational coefficients. Fix a prime $p$ and put $K=\Q(\zeta_p)$. The
argument includes $p=2$, although the principal application has $p$ odd.
Let $C\to\mathbf P^1$ be a smooth connected cyclic degree-$p$ cover with
chosen deck generator $\sigma$ and distinct branch points $t_1,\ldots,t_s$.
Write its nonzero residues as $a_i\in\F_p^\times$, with
$\sum_i a_i=0$. Assume $s\geq3$, and set $h=s-2$ and $J=J(C)$.
Then $\dim_K H^1(C,\Q)=h$ and the genus is $(p-1)h/2$.
The action of $K$ sends $\zeta_p$ to $\sigma^*$; use the compatible
action on $J$ under the Abel--Jacobi identification of $H^1$.

The rational determinant realization $D_C$ is specified by
\[
 (D_C)_\C=\bigoplus_{b=1}^{p-1}
       \bigwedge^h H^1(C,\C)_b,
 \qquad \dim_\Q D_C=p-1.
\]
It has two natural realizations: in $H^h(C^h,\Q)$ by antisymmetrized
K\"unneth tensors, and in
\[
 H^h(J,\Q)=\bigwedge^h H^1(J,\Q).
\]
An algebraic correspondence between these realizations is explained below.
Write
\[
 F_p^h=\{x_1^p+\cdots+x_{h+2}^p=0\}\subset\mathbf P^{h+1}.
\]

\begin{theorem}\label{thm:transfer}
There is a rational algebraic projector on $F_p^h$ whose cohomological
image, concentrated in degree $h$, is isomorphic to $D_C$ by algebraic
correspondences in both directions. For the realization in $C^h$, the
projector and the two inverse correspondences already satisfy the required
identities in rational Chow groups. In particular the realization of $D_C$
in $H^h(J,\Q)$ is the image of a rational algebraic correspondence from
$H^h(F_p^h,\Q)$.
\end{theorem}

\begin{corollary}\label{cor:products}
Let $C_1,\ldots,C_N$ be any finite list of such covers of the same prime
degree, with arbitrary branch points, residues and repetitions. Put
$h_j=s_j-2$ and $H=\sum_jh_j$. Every rational Hodge class in
\[
 \bigotimes_{\Q,j}D_{C_j}
 \ \subset\ H^H\left(\prod_jJ(C_j),\Q\right)
\]
is algebraic. The same holds in any rational algebraically projected
subspace of this tensor product, including the subspace with all
cyclotomic embedding labels equal.
\end{corollary}

No monodromy, genericity or endomorphism-algebra assumption appears here.
The assertion is about determinant subspaces, not the entire cohomology
of an arbitrary Jacobian or its powers.

\section{A common normal finite quotient}
Choose a coordinate on $\mathbf P^1$ for which all $t_i$ are finite.
With representatives $1\leq a_i\leq p-1$, the cover has a Kummer equation
\[
 y^p=\prod_{i=1}^{h+2}(x-t_i)^{a_i}.
\]
The point at infinity is unramified because $\sum_i a_i=pk$ for an integer
$k$. This equation describes the function field; $C$ denotes its smooth
projective normalization.

Let $Z=C^h$. On its $h$ factors act $\mu_p^h$, and let
\[
 N=\{(\epsilon_1,\ldots,\epsilon_h)\in\mu_p^h:
             \epsilon_1\cdots\epsilon_h=1\},\qquad
 G=N\rtimes\mathfrak S_h.
\]
Permutations permute the factors and normalize $N$. Let $Y=Z/G$, the
normal projective finite quotient. Its degree over
$\Sym^h\mathbf P^1=\mathbf P^h$ is $p$.

On the dense affine locus write
$u(T)=\prod_{\nu=1}^h(T-x_\nu)$. Its coefficients are coordinates on
$\Sym^h\mathbf A^1$. The $N$-invariant element
$z=\prod_{\nu=1}^h y_\nu$ is permutation invariant, and, after rescaling
$z$ by a constant, satisfies
\begin{equation}\label{eq:kummer}
 z^p=\prod_{i=1}^{h+2}u(t_i)^{a_i}.
\end{equation}
The coordinate-field extension of the ordered product has group
$\mu_p^h$ over $\C(x_1,\ldots,x_h)$: the $h$ Kummer classes are independent,
as seen by valuation at $x_\nu=t_i$. Thus taking $N$-invariants leaves
exactly the field generated by $z$. Taking permutation invariants gives
exactly \eqref{eq:kummer} over the symmetric coefficient field. This
also follows by comparing the degree $p$ of that equation with the index
of $G$ in $\mu_p^h\rtimes\mathfrak S_h$.

Homogeneously, evaluation at $t_i$ is a linear form $L_i$ on
$H^0(\mathbf P^1,\mathcal O(h))$. The $h+2$ hyperplanes $L_i=0$ are in
general position: any $h+1$ evaluations at distinct points are independent
by polynomial interpolation. Consequently there is a unique linear
relation, up to scaling,
\[
 \sum_{i=1}^{h+2}\lambda_i L_i=0,
 \qquad \lambda_i\ne0\quad\hbox{for every }i.
\]
The nonvanishing follows because a relation omitting one evaluation would
contradict that independence.

Consider the smooth diagonal hypersurface
\[
 X=\{\textstyle\sum_i\lambda_i x_i^p=0\}\subset\mathbf P^{h+1}.
\]
Over $\C$ it is isomorphic to $F_p^h$ by diagonal rescaling. The map
$[x_i]\mapsto[x_i^p]$ is a finite quotient onto the hyperplane
$\sum_i\lambda_i u_i=0$, identified with $\mathbf P^h$ by $u_i=L_i$.
Its group is
\[
 A=\mu_p^{h+2}/\mu_p^{\mathrm{diag}}.
\]
Because $\sum_i a_i=0$ modulo $p$, the formula
$\chi_a((\epsilon_i))=\prod_i\epsilon_i^{a_i}$ is a well-defined
nontrivial character of $A$. Put $B=\ker\chi_a$; then $|B|=p^h$.
The quotient $X/B$ has precisely the Kummer function field
\eqref{eq:kummer}. For example, in the chart $x_{h+2}\ne0$, the
semi-invariant function
\[
 w=\frac{\prod_i x_i^{a_i}}{x_{h+2}^{pk}}
 \quad\hbox{satisfies}\quad
 w^p=\frac{\prod_i u_i^{a_i}}{u_{h+2}^{pk}}.
\]
Since both quotients are normal finite covers of the same $\mathbf P^h$
with the same function field, they are isomorphic to its normalization in
that field. We have therefore constructed finite quotient maps
\begin{equation}\label{eq:quotients}
 r:X\longrightarrow Y\longleftarrow Z:q,
 \qquad d_X=\deg r=p^h,\quad d_Z=\deg q=p^{h-1}h!.
\end{equation}
The intermediate $Y$ may be singular. This causes no difficulty for the
next argument, which uses only cycles on smooth $X$ and $Z$.

\section{An elementary finite-correspondence identity}
For smooth projective varieties of equal dimension $h$, a correspondence
in $\CH^h(X\times Z)$ acts from the cohomology of $X$ to that of $Z$;
write $T^{\mathsf t}$ for its transpose. Let $[X\times_Y Z]$ mean the
dimension-$h$ fundamental cycle of the scheme-theoretic fiber product,
with its generic multiplicities.

\begin{lemma}\label{lem:finite}
In the situation \eqref{eq:quotients}, set $T=[X\times_Y Z]$. Then
\begin{align*}
 T\circ T^{\mathsf t}&=d_X\sum_{g\in G}[\Gamma_g]
                        &&\text{in }\CH^h(Z\times Z),\\
 T^{\mathsf t}\circ T&=d_Z\sum_{b\in B}[\Gamma_b]
                        &&\text{in }\CH^h(X\times X).
\end{align*}
\end{lemma}
\begin{proof}
The intersection support in each composition is finite over $Y$. In particular the
two pulled-back correspondences intersect properly in $Z\times X\times Z$
(or $X\times Z\times X$): the intersection has dimension $h$.
On a dense open subset of $Y$ both quotient maps are finite \'etale torsors.
There the composition from $Z$ to itself consists of each graph of $g\in G$
with multiplicity $d_X$: for a point and its $g$-translate one can choose
any of the $d_X$ points of $X$ above their common image. The other identity
is identical with $X$ and $Z$ interchanged.

These generic computations determine the entire dimension-$h$ cycles.
Indeed, a component supported above the excluded proper closed subset of
$Y$ has dimension at most $h-1$, since all the relevant maps are finite.
It cannot contribute a dimension-$h$ component to the proper intersection
or its pushforward. Thus the generic cycle identities hold globally, and
therefore also in the indicated Chow groups. No pushforward to singular
cohomology or resolution theorem has been used.
\end{proof}

\section{The determinant projector and transfer}
For a finite group $L$ acting on $Z$, write
$e_L=|L|^{-1}\sum_{\ell\in L}[\Gamma_\ell]$. Let
$\widetilde G=\mu_p^h\rtimes\mathfrak S_h$, containing $G$ as a normal
subgroup of index $p$, and set
\begin{equation}\label{eq:P}
 P=e_G-e_{\widetilde G}\in\CH^h(Z\times Z).
\end{equation}
Nested averaging gives $e_Ge_{\widetilde G}=e_{\widetilde G}e_G
=e_{\widetilde G}$, hence $P^2=P$. This formula uses the quotient deck
action through one slot; it never replaces it by simultaneous action on
all slots. In particular it remains valid when $p\mid h$.

\begin{lemma}\label{lem:cohom}
The only nonzero cohomological image of $P$ is in degree $h$, where
\[
 \im P=\bigoplus_{b=1}^{p-1}\bigwedge^hH^1(C,\C)_b
\]
after extension of scalars to $\C$.
\end{lemma}
\begin{proof}
K\"unneth tensors invariant under $N$ have equal deck-character labels
on their $h$ slots, because the characters trivial on $N$ are exactly
$(b,\ldots,b)$. The subtraction of $e_{\widetilde G}$ removes the label
$b=0$. A nontrivial character occurs only in $H^1(C)$, because the deck
action on $H^0(C)$ and $H^2(C)$ is trivial. Thus every surviving tensor
has total degree $h$ and lies in $H^1(C)_b^{\otimes h}$.

A geometric transposition of two factors acts on such K\"unneth tensors
as minus the ordinary tensor transposition, by graded commutativity.
Consequently taking geometric $\mathfrak S_h$-invariants takes ordinary
alternating tensors, namely $\bigwedge^hH^1(C)_b$. Each is a line since
$\dim H^1(C)_b=h$. The rationality is automatic from \eqref{eq:P}.
\end{proof}

Put $d=d_Xd_Z$ and define
\[
 R=\frac1d T^{\mathsf t}PT\in\CH^h(X\times X),\qquad
 a=PT,\qquad b=\frac1d T^{\mathsf t}P.
\]
Lemma~\ref{lem:finite} and $Pe_G=P$ give
\[
 R^2=R,\qquad ab=P,\qquad ba=R.
\]
For instance $PTT^{\mathsf t}P=dP$. These are Chow identities.
They exhibit $(Z,P)$ as a direct summand of $X$ with inverse algebraic
correspondences. Lemma~\ref{lem:cohom} proves the first realization of
Theorem~\ref{thm:transfer}.

\paragraph{Passage to the Jacobian.}
Let $\iota:C\to J$ be an Abel--Jacobi map based at a branch point; it is
equivariant for the stated deck convention and induces an isomorphism
$\iota^*:H^1(J,\Q)\to H^1(C,\Q)$. Its inverse is algebraic in degree one.
Here is a direct way to make this familiar fact sufficient for our use.
The inverse Hodge map, using the intersection pairing on $H^1(C)$, is a
rational Hodge class in the K\"unneth component
$H^1(C,\Q)\otimes H^1(J,\Q)\subset H^2(C\times J,\Q)$.
By the Lefschetz $(1,1)$ theorem it is the class of a rational divisor.
Composing this divisor correspondence with the curve projector
\[
 \pi_1^C=[\Delta_C]-[c_0\times C]-[C\times c_0]
\]
ensures its relevant source is precisely $H^1(C)$.
Call the resulting correspondence $L:C\rightsquigarrow J$.

Take $L^{\otimes h}$ from $C^h$ to $J^h$ and follow it by pullback along
the diagonal $J\to J^h$. On the alternating subspace this is the usual
wedge map. With the alternating inclusion normalized by $1/h!$, it
identifies the determinant realization in $C^h$ with that in $J$.
All maps are algebraic. For an explicit inverse on the determinant,
pull back by the sum map
$C^h\to J$, $(c_1,\ldots,c_h)\mapsto\sum_\nu\iota(c_\nu)$,
apply $(\pi_1^C)^{\otimes h}$, and divide by $h!$.
Its degree-$h$ component is the normalized alternating inclusion followed
by $(\iota^*)^{\otimes h}$. This proves the
remaining assertion of Theorem~\ref{thm:transfer}. No Hodge conjecture in
higher codimension is used in this degree-one passage.

\section{Products, rational Hodge classes, and balance}
We use the following exact published input.

\begin{theorem}[Shioda]\label{thm:shioda}
For a fixed prime $p$, the rational Hodge conjecture holds for every finite
product $\prod_j F_p^{h_j}$, with arbitrary dimensions $h_j$.
\end{theorem}

This is Theorem~2 on printed page 112 of \cite{Shioda79}; the paragraph
immediately preceding it verifies the condition used there for prime
degree and every dimension. The degree is fixed, but the dimensions need
not agree. This is stronger than the single-Fermat-hypersurface assertion.

To prove Corollary~\ref{cor:products}, first use the inverse degree-one
correspondences to place a rational Hodge class in
$\bigotimes_j\im P_j\subset H^H(\prod_j C_j^{h_j},\Q)$.
The exterior product of the $b_j$ constructed above maps it to a rational
Hodge class on $\prod_j F_p^{h_j}$; the exterior product of the $a_j$
returns the original class. If $H$ is even, Shioda makes the source
class algebraic, and its image under the forward algebraic correspondence
is algebraic. If $H$ is odd there is no rational $(r,r)$ class in degree
$H$. This proves the corollary. The explicit algebraic splitting means
that even semisimplicity of polarizable Hodge structures is unnecessary
for this proof.

For comparison with the companion fusion statement, the matched space is
\[
 W_\C=\bigoplus_{b=1}^{p-1}
             \bigotimes_j\bigwedge^{h_j}H^1(C_j,\C)_b.
\]
It is rational. To realize matching by an algebraic projector, one may
work in all $H=\sum h_j$ curve slots and use
\[
 \prod_{r=2}^{H}\left(\frac1p\sum_{u=0}^{p-1}
                  (\sigma_1^*)^u(\sigma_r^*)^{-u}\right).
\]
On the already alternating factors it selects exactly equal labels in
all slots. This formulation also works when one or several $h_j$ are
divisible by $p$; simultaneous action on $h_j$ slots is not substituted
for the action on one selected slot.

Write $[a]_p\in\{1,\ldots,p-1\}$ for the representative of a nonzero
residue, and put $S=\sum_js_j$. The eigenspace Hodge numbers of a cyclic
cover give the criterion
\begin{equation}\label{eq:balance}
 2\sum_{j,i}[b a_{ji}]_p=pS\qquad (b\in\F_p^\times)
\end{equation}
for every component line of $W$ to have bidegree $(H/2,H/2)$.
Changing the pullback labeling interchanges $b$ with $-b$, leaving the
criterion unchanged. Hence \eqref{eq:balance} and
Corollary~\ref{cor:products} imply
\[
 W\subset\im\left(\cl:\CH^{H/2}\left(\prod_jJ(C_j)\right)_\Q
                       \longrightarrow H^H\left(\prod_jJ(C_j),\Q\right)\right).
\]
There is no need to decompose a balanced residue tuple into opposite pairs,
construct a degeneration, or assume that an individual determinant is a
Hodge class. Jointly balanced products of non-Hodge determinants are
included directly.

\begin{remark}[A previously proposed test example]
For $p=5$, the branch tuples $(1,1,3)$ and $(1,2,4,4,4)$ give $h_1=1$
and $h_2=3$. Their residue-age vectors, for multipliers $1,2,3,4$, are
$(1,1,2,2)$ and $(3,3,2,2)$. Their sums are all four, so the matched
rational four-dimensional space in
$H^1(J(C_1))\otimes H^3(J(C_2))$ has type $(2,2)$.
Its algebraicity is a consequence of the transfer from
$F_5^1\times F_5^3$ and Shioda's theorem. The fact that neither individual
determinant is of Tate type does not distinguish this example from the
Fermat-product method.
\end{remark}

\section{Attribution and limits}
The cycle theorem doing the essential algebraicity work is Shioda's
published Fermat-product theorem \cite{Shioda79}. Schoen's 1989 paper
\cite{Schoen89}, in its directly inspected introduction on printed pages
137--138, announces algebraic correspondences from cyclic covers of
$\mathbf P^v$ branched over $v+2$ hyperplanes, which are Fermat quotients,
to determinant Hodge structures of cyclic-cover generalized Pryms.
That is direct precedent for the construction used here, and must be
acknowledged prominently.

At the time of writing, only those two introduction pages of
\cite{Schoen89} were available to this audit; the exact hypotheses of its
later sections were not inspected. Accordingly we do not cite an
unverified numbered theorem as if it already had precisely the scope of
Theorem~\ref{thm:transfer}. Instead Sections 2--4 give the needed
$\mathbf P^1$ argument explicitly. This access limitation does not become
a claim of novelty. The product consequence is a short deduction from
classical Fermat theory once this finite-quotient correspondence is written
down; identifying its earliest explicit publication remains a bibliographic
task.

The companion all-powers assertion for a very general elementary-abelian
cover additionally needs the simultaneous monodromy classification and
the algebraic invariant-generator reduction. Those assertions do not
follow merely from this note. Conversely the arbitrary balanced determinant
algebraicity step should not be described as a new cycle-construction
breakthrough on the evidence presently available.

\begin{thebibliography}{9}
\bibitem{Shioda79}
T.~Shioda, \emph{The Hodge conjecture and the Tate conjecture for Fermat
varieties}, Proc. Japan Acad. Ser. A Math. Sci. \textbf{55} (1979), 111--114,
Theorem~2, p.~112.
\href{https://doi.org/10.3792/pjaa.55.111}{doi:10.3792/pjaa.55.111}.
\href{https://www.researchgate.net/profile/Tetsuji-Shioda/publication/226509160_The_hodge_conjecture_for_Fermat_varieties/links/56716efe08ae2b1f87aef843/The-hodge-conjecture-for-Fermat-varieties.pdf}{Author-uploaded original scan (direct PDF)}.
The landing-page metadata labels the later Math. Ann. paper; the uploaded
four-page primary scan is the Proc. Japan Acad. note cited here.

\bibitem{Schoen89}
C.~Schoen, \emph{Cyclic covers of $\mathbf P^v$ branched along $v+2$
hyperplanes and the generalized Hodge conjecture for certain abelian
varieties}, in Arithmetic of Complex Manifolds, Lecture Notes in Math.
\textbf{1399}, Springer, 1989, 137--154.
\href{https://doi.org/10.1007/BFb0095974}{doi:10.1007/BFb0095974}.
\href{https://page-one.springer.com/pdf/preview/10.1007/BFb0095974}{Publisher introduction preview}.

\bibitem{Schoen88}
C.~Schoen, \emph{Hodge classes on self-products of a variety with an
automorphism}, Compositio Math. \textbf{65} (1988), 3--32.
\href{https://www.numdam.org/item/CM_1988__65_1_3_0.pdf}{Original journal scan}.

\bibitem{Aoki02}
N.~Aoki, \emph{Hodge cycles on CM abelian varieties of Fermat type},
Comment. Math. Univ. Sancti Pauli \textbf{51} (2002), 99--130,
Theorem~1.2, p.~102.
\href{https://rikkyo.repo.nii.ac.jp/record/8744/files/AA00610867_51-01_07.pdf}{Original university-hosted PDF}.
\end{thebibliography}
\end{document}
